
\فصل{روش پیشنهادی}

انطباق در زمان آزمون برخط یکی از مسائل کلیدی در یادگیری ماشین مدرن به شمار می‌رود. در بسیاری از کاربردهای دنیای واقعی، داده‌های ورودی در طول زمان دچار تغییر توزیع می‌شوند و این تغییر امری اجتناب‌ناپذیر است. از آنجا که داده‌های جدید معمولاً فاقد برچسب هستند، لازم است مدل بتواند به صورت پیوسته، برخط، در لحظه و به صورت بلادرنگ با داده‌های جدید سازگار شود. چنین قابلیتی امکان استفاده‌ی پایدار از مدل‌های یادگیری عمیق را در محیط‌های پویا فراهم می‌سازد. با این حال، اجرای این فرایند با چالش‌هایی همچون انباشت خطا و فراموشی فاجعه‌بار همراه است که می‌تواند موجب تضعیف عملکرد مدل در طول زمان گردد.

یکی از خلأهای مهم در روش‌های موجود، فقدان بررسی کافی در زمینه‌ی حساسیت و استحکام مدل نسبت به تغییرات فراپارامترها است. در شرایط خاص مسئله‌ی انطباق در زمان آزمون برخط که هیچ دانشی از توزیع هدف در دسترس نیست، عملا امکان تنظیم دقیق و هدفمند فراپارامترها وجود ندارد. در بسیاری از مسائل یادگیری با ناظر یا یادگیری بدون ناظر، اصطلاحا {دوستانه‌بودن فراپارامترها\پانویس{Hyper-parameter Friendliness}} یک مزیت و ترجیح محسوب می‌شود؛ اما در مسئله‌ی انطباق برخط در زمان آزمون، این موضوع یک ضرورت حیاتی است. به بیان دیگر، رویکرد مطلوب باید در برابر تغییرات فراپارامترها مقاوم بوده و به فرایند {تنظیم دقیق\پانویس{Fine-tune}} وابسته نباشد.

از سوی دیگر، روش‌های پیشین برای شناسایی و حذف نمونه‌های نامطمئن یا مخرب، عمدتا به معیارهایی نظیر آنتروپی یا بیشینه‌ی احتمالات خروجی شبکه متکی شده‌اند. این معیارها اساساً برای داده‌های هم‌توزیع طراحی شده‌اند و صحت عملکرد آن‌ها هنگام مواجهه با داده های زمان آزمون که لزوما توزیع مشابهی با داده‌های اولیه که مدل روی آن‌ها آموزش دیده ندارند، می تواند کاملا غیرقابل اعتماد باشد.

با توجه به چالش‌های مطرح شده در ادامه‌ی این فصل به دنبال ارائه ی روشی هستیم که علاوه بر سادگی در پیاده‌سازی، نسبت به تغییرات فراپارامترها حساسیت کمی داشته و استحکام بالاتری را ارائه دهد. همچنین، این روش سازوکاری جدید برای شناسایی نمونه‌های نامطمئن و خارج از توزیع فراهم می‌کند؛ سازوکاری که به‌طور ویژه متناسب با شرایط تغییر توزیع در مسئله‌ی انطباق برخط زمان آزمون طراحی شده است.

\قسمت{پدیده‌ی اطمینان بیش‌ازحد در شبکه‌های عصبی} 

یکی از مواردی که پژوهش‌های پیشین انطباق در زمان آزمون، به‌ویژه در حالت برخط، در نظر نگرفته‌اند، پدیده‌ی اطمینان بیش‌ازحد شبکه‌های عصبی در مواجهه با نمونه‌های خارج از توزیع است. اگرچه شبکه‌های عصبی عمیق به نتایج قابل توجهی در زمینه‌های مختلف یادگیری ماشین دست یافته‌اند، اما یکی از چالش‌های اساسی در این شبکه‌ها همین پدیده‌ی اطمینان بیش‌ازحد در مواجهه با داده‌های خارج از توزیع است. به بیان ساده، مدل نه‌تنها قادر به شناسایی داده‌های خارج از توزیع نیست، بلکه در بعضی از موارد با اطمینان بالایی پاسخ نادرست ارائه می‌دهد. این رفتار می‌تواند منجر به خطاهای جدی در کاربردهای حساس شود و به‌ویژه در مسئله‌ی انطباق زمان آزمون برخط، 
اهمیت ویژه‌ای دارد.  
علت اصلی این امر آن است که در انطباق زمان آزمون برخط، دسترسی به برچسب داده‌های ورودی در زمان آزمون و امکان دریافت بازخورد لحظه‌ای در بسیاری از موارد وجود ندارد
و فرآیند انطباق صرفاً بر اساس ساختار داده‌ها و پیش‌بینی‌های مدل انجام می‌شود. 
در چنین شرایطی، تکیه بر معیارهای درست‌نمایی\پانویس{Likelihood} یا خروجی‌های مستقیم مدل 
که لزوماً نشان‌دهنده‌ی میزان اعتمادپذیری پیش‌بینی نیستند، 
می‌تواند مدل را به‌سمت انطباق نادرست سوق دهد.  
بنابراین، نادیده گرفتن مسئله‌ی اطمینان بیش‌ازحد در طراحی الگوریتم‌های انطباق زمان آزمون، 
خطر شکست و فروپاشی مدل را در مواجهه با داده‌های غیرمنتظره و خارج از توزیع به‌طور جدی افزایش می‌دهد.  

از این‌رو، لازم است سازوکار‌هایی برای سنجش و کنترل سطح اطمینان مدل در فرآیند انطباق برخط در نظر گرفته شوند 
تا از تأثیرگذاری منفی داده‌های خارج از توزیع جلوگیری گردد. 
رویکرد پیشنهادی  نیز دقیقاً با در نظر گرفتن این شکاف طراحی شده است 
تا انطباق زمان آزمون برخط در حضور داده‌های خارج از توزیع با پایداری و دقت بیشتری انجام شود.
در ادامه، ابتدا این پدیده را با جزئیات بیشتری مورد بررسی قرار داده 
و سپس یک دسته‌بندی کلی از روش‌های موجود در حوزه‌ی شناسایی خارج از توزیع ارائه می‌کنیم 
تا جایگاه رویکرد پیشنهادی در میان آن‌ها روشن شود.


در همین راستا روش‌های \cite{class-ood1, class-ood2, class-ood3}
نشان می دهند که شبکه‌های عصبی قادرند برای نمونه‌هایی که به هیچ‌یک از کلاس‌های آموزشی تعلق ندارند، 
پیش‌بینی‌هایی با اطمینان بیش‌ازحد ارائه دهند. 
چنانچه سامانه قادر به شناسایی این ورودی‌ها به‌عنوان داده‌های ناشناخته نبوده و فرایند به‌روز‌رسانی پارامترهای شبکه را روی این نمونه‌ها نیز انجام دهد، 
پیامدهای خطرناکی به همراه خواهد داشت. از این رو، شناسایی داده‌های خارج از توزیع
برای تضمین ایمنی و قابلیت اعتماد در به‌کارگیری شبکه‌های عصبی امری حیاتی است \cite{ood1, ood2}.

هم چنین روش‌های \cite{Nalisnick2019, Choi2019, Shafaei2018} نشان داده‌اند که {مدل‌های مولد عمیق\پانویس{Deep Generative Models}} در شناسایی داده‌های آموزشی از داده‌های خارج از توزیع بر اساس معیار درست‌نمایی ناتوان هستند. این پدیده نه تنها زمانی رخ می‌دهد که داده‌ها شباهت ظاهری دارند، بلکه حتی در شرایطی که معنا و محتوای آن‌ها کاملاً متفاوت است نیز مشاهده شده است. به عنوان مثال، مدل \lr{Glow} \cite{Kingma2018} که یک جریان نرمال‌ساز پیشرفته محسوب می‌شود، وقتی بر روی مجموعه‌ی \lr{CIFAR-10} آموزش دیده باشد، برای داده‌های \lr{SVHN} درست‌نمایی بالاتری نسبت به داده‌های آموزشی \lr{CIFAR-10} اختصاص می‌دهد \cite{Nalisnick2019, Choi2019}. این نتیجه شگفت‌آور است؛ زیرا \lr{CIFAR-10} شامل تصاویر حیوانات و وسایل نقلیه است، در حالی که \lr{SVHN} شامل اعداد خیابانی می‌باشد. به طور طبیعی، یک انسان هرگز این دو مجموعه را با یکدیگر اشتباه نمی‌گیرد. 
از منظر الگوریتمی نیز این یافته‌ها نگران‌کننده هستند، زیرا منجر به شکست بسیاری از روش‌های پیشنهادی برای اعتبارسنجی دسته‌بندها \cite{Bishop1994} و همچنین شناسایی ناهنجاری \cite{Pimentel2014} می‌شوند.

به طور خاص به منظور ارائه‌ی یک نمونه راه حل برای مقابله با این مشکل، \cite{typicality} روشی مبتنی بر {اصول نمونه‌واری\پانویس{Typicality}} را معرفی می‌کند. این روش بیان می‌کند داده‌های متعلق به یک توزیع مشخص باید دارای مقادیر درست‌نمایی قرارگرفته در یک بازه‌ی مشخص باشند. ایده‌ی اصلی این روش بر پایه‌ی مقایسه‌ی درست‌نمایی یک دسته از نمونه‌ها با { درست‌نمایی مورد انتظار\پانویس{Expected Likelihood}} یعنی میانگین درست‌نمایی داده‌های آموزشی که به عنوان مرجع در نظر گرفته می‌شود است. اگر درست‌نمایی دسته‌ی داده‌های ورودی به طور معناداری بیشتر یا کمتر از مقدار مورد انتظار باشد، آن دسته به عنوان داده‌ی خارج از توزیع شناسایی می‌شود.

فرض می‌شود که $p_\theta(x)$ درست‌نمایی نمونه‌ی $x$ تحت مدل با پارامترهای $\theta$ باشد. مقدار مورد انتظار درست‌نمایی با استفاده از داده‌های آموزشی به صورت زیر طبق رابطه‌ی \ref{eq:expected-likelihood} تخمین زده می‌شود\cite{typicality}:
\begin{equation}
	\label{eq:expected-likelihood}
	\mathbb{E}_{x \sim \mathcal{D}_{\text{train}}}\left[ \log p_\theta(x) \right],
\end{equation}
که در آن $\mathcal{D}_{\text{train}}$ مجموعه‌ی داده‌های آموزشی است. سپس، برای یک دسته‌ی ورودی $\mathcal{B}$، میانگین درست‌نمایی طبق رابطه‌ی \ref{eq:batch-likelihood} محاسبه می‌شود\cite{typicality}:
\begin{equation}
	\label{eq:batch-likelihood}
	\hat{L}(\mathcal{B}) = \frac{1}{|\mathcal{B}|} \sum_{x \in \mathcal{B}} \log p_\theta(x).
\end{equation}

بر اساس معیار نمونه‌واری، اگر $\hat{L}(\mathcal{B})$ به‌طور معناداری خارج از بازه‌ی تعریف‌شده پیرامون مقدار \eqref{eq:expected-likelihood} قرار گیرد، دسته‌ی $\mathcal{B}$ به عنوان خارج از توزیع طبق رابطه‌ی \ref{eq:ood-detection} برچسب زده می‌شود\cite{typicality}:
\begin{equation}
	\label{eq:ood-detection}
	\text{OOD}(\mathcal{B}) =
	\begin{cases}
		1 & \text{if } \hat{L}(\mathcal{B}) \notin [\mu - \epsilon, \mu + \epsilon], \\
		0 & \text{otherwise}.
	\end{cases}
\end{equation}

که در آن $\mu$ مقدار مورد انتظار طبق رابطه‌ی \ref{eq:expected-likelihood} و $\epsilon$ آستانه‌ای کوچک برای تعیین بازه‌ی قابل‌قبول است. 

این روش نشان می‌دهد که صرفاً اتکا به درست‌نمایی خام در مدل‌های شبکه عصبی می‌تواند گمراه‌کننده باشد، و استفاده از اصل نمونه‌واری می‌تواند به عنوان ابزاری کارآمد برای شناسایی نمونه‌های خارج از توزیع مورد استفاده قرار گیرد.

\زیرزیرقسمت{شناسایی خارج از توزیع در مرحله‌ی آموزش}

همان‌طور که پیش‌تر بیان شد از آنجا که شبکه‌های عصبی بر اساس {فرض جهان بسته\پانویس{Closed-world}} طراحی می‌شوند \cite{krizhevsky2017imagenet}، 
با مشکل اطمینان بیش‌ازحد روی داده‌های خارج از توزیع مواجه هستند، 
یعنی به داده‌های خارج از توزیع نیز پیش‌بینی‌هایی با سطح اطمینان مشابه داده‌های درون‌توزیع اختصاص می‌دهند \cite{class-ood2}. 
این مسئله تهدیدی جدی برای کاربرد شبکه‌های عصبی در شرایط {جهان باز\پانویس{Open-world}} محسوب می‌شود.  

برای رفع این مشکل در مرحله‌ی آموزش یعنی هنگام کار روی هدف اصلی مسئله، روش‌های مختلفی پیشنهاد شده‌اند. 
برای نمونه، \lr{OE} \cite{hendrycks2018deep} 
از مجموعه‌داده‌های کمکی خارج از توزیع برای واداشتن مدل به پیش‌بینی با اطمینان پایین‌تر روی آن‌ها بهره می‌گیرد 
و بدین ترتیب شبکه را بهتر {واسنجی\پانویس{Calibrate}} می‌کند.  
روش‌های جدیدتر مانند \lr{DOE}\cite{wang2023out} و \lr{MixOE}\cite{zhang2023mixture} این رویکرد را توسعه داده‌اند. 
از سوی دیگر، \cite{tang2023codes} با بهره‌گیری از مدل‌های مولد مانند \lr{Chamfer GAN} داده‌های خارج از توزیع مصنوعی ایجاد کردند.  
روش‌هایی نظیر \cite{du2022vos} و \cite{tao2023npos} نیز داده‌های خارج از توزیع مجازی را در سطح غیرتصویری تولید کرده 
و بدون نیاز به تصاویر واقعی خارج از توزیع، عمومیت ‌بخشی به \lr{OE} را تقویت می‌کنند.  

اگرچه این روش‌ها نتایج چشمگیری به دست آورده‌اند، 
اما همگی مستلزم بازآموزی کامل شبکه و استفاده از مجموعه داده‌های کمکی هستند. 
وجود این وابستگی علاوه بر افزایش پیچیدگی، می‌تواند دقت شبکه را روی مسئله‌ی اصلی کاهش دهد و محدودیت‌های خاصی را در انعطاف‌پذیری مدل ایجاد نماید. 
از آنجا که این دسته از روش‌ها باید در زمان طراحی و آموزش اولیه‌ی مدل در نظر گرفته شوند، 
در عمل در دسته‌ی مشابه مسائل {آموزش در زمان آزمون\پانویس{Test-time Training}} \cite{sun2020ttt, gandelsman2022tttmae} قرار می‌گیرند که نیاز‌مند طراحی یک {مسئله‌ی کمکی\پانویس{Auxiliary Task}} در کنار مسئله‌ی اصلی هستند.

\زیرزیرقسمت{شناسایی خارج از توزیع پسینی}

تمرکز مهم دیگر در پژوهش‌های اخیر، توسعه‌ی روش‌های {شناسایی خارج از توزیع پسینی\پانویس{Post-hoc OOD Detection}} است که بدون نیاز به بازآموزی شبکه عمل می‌کنند.
برای نمونه، \cite{hendrycks2016baseline} {الگوریتم بیشینه‌ی احتمال هموار\پانویس{Maximum Softmax Probability}} را معرفی کردند
که فرض می‌کند داده‌های خارج از توزیع در مقایسه با داده‌های درون‌توزیع سطح اطمینان کمتری خواهند داشت. 
\cite{liang2018enhancing} روش \lr{ODIN}\پانویس{Out-of-DIstribution detector for Neural networks} را با استفاده از {مقیاس‌گذاری دما\پانویس{Temperature Scaling}} ارائه می‌کند
تا شکاف میان داده‌های درون توزیع و خارج از توزیع افزایش یابد.  
\cite{liu2020energy} امتیاز انرژی را برای کاهش سوگیری در مدل‌های مبتنی بر انرژی به کار می‌گیرد.  
\cite{sun2022dice} روش \lr{DICE} \پانویس{Directed Sparisification}  را ارائه می‌کند که بر اساس انتخاب {نورون‌های\پانویس{Neuron}} مهم عمل می‌کند 
هم‌چنین روش \lr{ReAct} \cite{sun2021react} \پانویس{Rectified Activations} خروجی‌ها را با تصحیح فعال‌سازی برای شناسایی بهتر نمونه‌های درون توزیع و خارج از توزیع پالایش می‌کند.  
\cite{yuan2023lhact} این رویکرد را با معرفی \lr{LHAct} \پانویس{Low and High Activations} گسترش داده و 
با استفاده از فیلتر {باترورث‌\پانویس{Butterworth}} مقید به اصلاح فعال‌سازی‌ها شبکه پرداختند.  
\cite{tang2024cores} نیز با ردیابی لایه‌های {تماما متصل\پانویس{Fully-connected}} و طراحی امتیازهای خارج از توزیع بر مبنای هسته‌های مرتبط با نمونه، 
راهکاری نوین ارائه کردند.  
همچنین \cite{zhu2022boosting} ویژگی‌ها را به یک مجموعه‌ی نمونه‌وار تصحیح کرده و 
سازگاری روش خود را با الگوریتم‌های مختلف شناسایی خارج از توزیع نشان می‌دهد. 

با وجود پیشرفت‌های چشمگیر، بیشتر این روش‌ها از امتیازهای مشتق‌شده از وظیفه‌ی دسته‌بندی داده‌های درون‌توزیع بهره می‌گیرند 
و برای شناسایی خارج از توزیع به‌طور خاص طراحی نشده‌اند. 
در نتیجه، ممکن است در شناسایی داده‌های خارج از توزیع با محدودیت‌هایی مواجه شوند و راه‌حل بهینه‌ای را ارائه نکنند.\\
روش پیشنهادی در دسته‌ی روش‌های پسینی قرار می‌گیرد. ماهیت مسائل انطباق برخط در زمان آزمون مبتنی بر مدل‌های از پیش آموزش‌دیده است؛ در بسیاری از موارد دسترسی به مدل خام اولیه و به‌ویژه مجموعه‌داده‌ی آموزشی امکان‌پذیر نیست. حتی در صورت دسترسی، محدودیت‌های منابع محاسباتی، ملاحظات حریم خصوصی و زمان‌بر بودن بازآموزی، استفاده‌ی مجدد از داده‌های اولیه را عملا ناممکن می‌سازد. همین چالش‌ها زمینه‌ساز توسعه‌ی روش‌های انطباق زمان آزمون شده‌اند. بر این اساس، جایگاه رویکرد روش پیشنهادی در این حوزه بهره‌گیری از روش‌های شناسایی خارج از توزیع پسینی است که در بخش بعدی به‌طور مفصل بررسی خواهد شد.

\قسمت{معیار جدید: امتیاز جرم تقریبی}

مطالعات اخیر نشان داده‌اند که استفاده‌ی صرف از درست‌نمایی برای شناسایی داده‌های خارج از توزیع در ابعاد بالا کافی نیست \cite{Nalisnick2019}. 
به بیان دیگر، ممکن است یک نمونه دارای درست‌نمایی بالا تحت یک توزیع باشد، 
اما در عمل احتمال مشاهده‌ی آن بسیار ناچیز باشد. 
این مسئله به‌ویژه در فضاهای با ابعاد بالا تشدید می‌شود، 
زیرا در چنین شرایطی اغلب جرم احتمالی به‌صورت متمرکز در ناحیه‌ای محدود از فضای ورودی توزیع می‌شود 
و نمونه‌هایی که حتی اندکی از این ناحیه فاصله دارند، 
به‌رغم داشتن مقدار درست‌نمایی قابل توجه، در عمل بسیار بعید هستند.  
این موضوع در حالت انطباق در زمان آزمون اهمیت بیشتری می‌یابد، 
چرا که در این شرایط مدل ناگزیر است بدون دسترسی به داده‌های آموزشی اولیه 
و صرفاً بر اساس نمونه‌های ورودی جاری تصمیم‌گیری کند؛
بنابراین معیارهای حساس به موقعیت نمونه در {مجموعه‌ی نمونه‌وار\پانویس{Typical Set}} 
می‌توانند به مدل کمک کنند تا رفتار خود را به‌طور هوشمندانه تطبیق دهد.  

نمونه‌های واقعی یک توزیع معمولاً در مجموعه‌ای موسوم به مجموعه‌ی نمونه‌وار قرار دارند. 
این مجموعه به ناحیه‌ای اطلاق می‌شود که بخش عمده‌ی جرم احتمالی در آن متمرکز است 
و رفتار آماری داده‌ها در آن به‌خوبی با توزیع اصلی مطابقت دارد.
این مسئله پیش از این در شبکه‌های مولد عمیق مورد بررسی قرار گرفته‌است \cite{typicality}.
اگرچه یک نقطه می‌تواند در ناحیه‌ای با چگالی بالا قرار گیرد، 
اما در صورتی‌که پیرامون آن ناحیه از چگالی بسیار کمی برخوردار باشد، 
آن نقطه عملاً در خارج از مجموعه‌ی نمونه‌وار قرار دارد 
و نمی‌توان آن را به‌عنوان نمونه‌ای واقعی از توزیع در نظر گرفت.  
این تمایز میان چگالی نقطه‌ای و جرم ناحیه‌ای، هسته‌ی اصلی مسئله در شناسایی داده‌های خارج از توزیع محسوب می‌شود 
و در محیط انطباقی برخط، می‌تواند مبنایی برای تصمیم‌گیری پویا درباره‌ی این‌که کدام داده‌ها 
در فرآیند به‌روزرسانی مدل مؤثر واقع شوند، فراهم آورد.  

بر این اساس، انتظار می‌رود در نقاطی با درست‌نمایی بالا که خارج از مجموعه‌ی نمونه‌وار قرار دارند، 
چگالی به‌سرعت تغییر کند. 
به‌عبارت دیگر، در چنین نقاطی منحنی توزیع به‌گونه‌ای است که حتی یک جابه‌جایی بسیار کوچک در فضای ورودی منجر به تغییر شدید مقدار درست‌نمایی می‌شود. 
این ویژگی به‌طور مستقیم در اندازه‌ی گرادیان لگاریتم درست‌نمایی نسبت به ورودی منعکس می‌شود.  
اگر $\nabla_x \log p_\theta(x)$ را گرادیان لگاریتم درست‌نمایی نسبت به ورودی در نظر بگیریم، 
میزان بزرگی این گرادیان نشان‌دهنده‌ی شدت تغییرات چگالی در ناحیه‌ی پیرامونی نمونه است.  
به‌گونه‌ای که اندازه‌ی این گرادیان برای داده‌های خارج از توزیع 
به‌طور معناداری بزرگ‌تر از نمونه‌های درون‌توزیع خواهد بود، 
چراکه آن‌ها در نواحی با جرم احتمالی کم قرار دارند.  
این خاصیت برای مسئله‌ کلیدی است: 
چون در فرایند انطباق زمان آزمون لازم است نمونه‌های خارج از توزیع شناسایی شده 
و سهم آن‌ها در به‌روزرسانی یا تنظیم پارامترهای مدل به شکل کنترل‌شده مدیریت گردد.  

از این‌رو، معیار جدیدی برای امتیازدهی به داده‌ها به صورت زیر طبق رابطه‌ی \ref{eq:apms} تعریف می‌شود:
\begin{equation}
	s_\theta(x) = - \left\| \nabla_x H(f_{\bm\theta}(x)) \right\|^2
	\label{eq:apms}
\end{equation}

این امتیاز که با نام {امتیاز جرم تقریبی\پانویس{Approximate Mass Score}} شناخته می‌شود، 
قابلیتی جدید برای شناسایی داده‌های خارج از توزیع فراهم می‌سازد. 
برتری این معیار در آن است که به جای اتکا به چگالی نقطه‌ای، 
اطلاعات محلی مربوط به تغییرات چگالی در اطراف هر نمونه را نیز در نظر می‌گیرد. 
به بیان دیگر، این امتیاز نوعی برآورد تقریبی از جرم احتمالی ناحیه‌ی اطراف نمونه ارائه می‌کند 
و بدین ترتیب قادر است نمونه‌هایی را که به‌رغم داشتن درست‌نمایی بالا خارج از مجموعه‌ی نمونه‌وار قرار دارند، شناسایی نماید.  

از منظر انطباق در زمان آزمون، امتیاز جرم تقریبی این امکان را فراهم می‌آورد که 
مدل در مواجهه با جریان داده‌های ورودی تصمیم بگیرد کدام نمونه‌ها شایستگی مشارکت در انطباق را دارند 
و کدام داده‌ها به دلیل احتمال بالای خارج از توزیع بودن باید حذف شده یا وزن کمتری بگیرند.  
این امر نه‌تنها دقت مدل را در مواجهه با داده‌های ناشناخته ارتقا می‌دهد، 
بلکه از پدیده‌ی انباشت خطا و فراموشی فاجعه‌بار در حین به‌روزرسانی وزن‌های شبکه نیز جلوگیری می‌کند.  

به‌طور خلاصه، امتیاز جرم تقریبی یک ابزار مکمل برای روش‌های انطباق در زمان آزمون به شمار می‌رود 
که با ترکیب اطلاعات نقطه‌ای و ناحیه‌ای، 
نه‌تنها در شناسایی داده‌های خارج از توزیع مؤثر است 
بلکه در مدیریت فرآیند به‌روزرسانی مدل نیز نقش کلیدی ایفا می‌کند. 

\قسمت{انطباق در زمان آزمون با استفاده از امتیاز جرم تقریبی}

در بخش قبل، معیار جدیدی موسوم به امتیاز جرم تقریبی معرفی شد که بر مبنای اندازه‌ی گرادیان لگاریتم درست‌نمایی تعریف می‌شود 
و توانایی بالایی در تمایز میان داده‌های درون‌توزیع و خارج از توزیع از خود نشان می‌دهد. 
یادآوری می‌کنیم که صرف درست‌نمایی نقطه‌ای برای سنجش کیفیت نمونه کافی نیست، 
چرا که یک نقطه می‌تواند دارای مقدار درست‌نمایی بالا باشد، 
اما در عمل به دلیل قرار گرفتن در نواحی کم‌جرم از فضای ورودی، خارج از مجموعه‌ی نمونه‌وار محسوب شود. 
امتیاز جرم تقریبی با لحاظ کردن تغییرات محلی چگالی، 
اطمینان مناسب‌تری از تعلق یک نمونه به توزیع اصلی ارائه می‌دهد 
و همین ویژگی، آن را به معیاری کارآمد برای هدایت فرآیند انطباق در زمان آزمون تبدیل می‌کند.  
این امتیاز به مدل امکان می‌دهد با درک دقیق‌تر ساختار توزیع داده‌ها، تصمیم بگیرد کدام نمونه‌ها برای بروزرسانی وزن‌ها ارزشمند هستند و کدام داده‌ها می‌توانند موجب بی‌ثباتی شوند. به بیان دیگر، این معیار یک ابزار هوشمندانه برای تمایز نمونه‌های مفید و نامطمئن در جریان داده‌های برخط فراهم می‌آورد.

در حالت انطباق در زمان آزمون، با یک جریان پیوسته از داده‌های ورودی مواجه هستیم 
که همگی لزوماً کیفیت یکسانی ندارند. 
برخی از این داده‌ها ذاتاً دارای درست‌نمایی کم و همچنین جرم احتمالی ناچیز در اطراف خود هستند، 
به‌گونه‌ای که ورود مستقیم آن‌ها به فرآیند انطباق می‌تواند منجر به انباشت خطا و فروپاشی تدریجی مدل شود. 
برای مقابله با این مسئله، ابتدا یک مرحله‌ی پالایش نمونه‌ها به صورت زیر در نظر می‌گیریم.
نمونه‌هایی که آنتروپی پیش‌بینی آن‌ها بالاتر از یک آستانه‌ی مشخص $E_0$ باشد، حذف می‌شوند.  
در اینجا این فیلتر را با یک {تابع شاخص\پانویس{Identity-like Function}} طبق رابطه‌ی \ref{eq:hentfiltering} تعریف می‌کنیم:
\begin{equation}
	\label{eq:hentfiltering}
	S_{\mathrm{ent}}(x) \coloneqq \mathbb{I} \{ E(x; \Theta) < E_0 \},
\end{equation}
که $\mathbb{I}\{\cdot\}$ تابع شاخص است و تنها نمونه‌های با آنتروپی کمتر از $E_0$ را عبور می‌دهد.  
این فیلتر موجب می‌شود که داده‌های آشکارا غیرقابل اعتماد در همان ابتدا از فرایند انطباق کنار گذاشته شوند 
و تنها نمونه‌های با اطمینان نسبی وارد فرآیند وزن‌دهی و به‌روزرسانی شوند.  

در گام دوم، یک سازوکار وزن‌دهی مبتنی بر امتیاز جرم تقریبی تعریف می‌کنیم. 
ایده‌ی اصلی آن است که نمونه‌های با جرم تقریبی بالاتر یعنی داده‌هایی که علاوه بر درست‌نمایی مناسب در نواحی پرجرم نیز قرار دارند،
در فرآیند انطباق وزن بیشتری دریافت کنند. 
به‌طور مشخص، در هر لحظه از فرایند انطباق در زمان آزمون برخط یک دسته‌ از داده‌های جاری $X_t=\{x_1,\dots,x_m\}$ 
از جریان ورودی استخراج می‌شود. 
از این دسته هم‌زمان برای پیش‌بینی و انطباق استفاده می‌شود.
پس از انجام پیش‌بینی از این دسته برای انطباق استفاده می‌شود، ابتدا نمونه‌هایی که آنتروپی آن ها از یک آستانه‌ی مشخص $E_0$ بالاتر باشد از فرایند انطباق حذف می‌شوند و داریم:  $X'_t = \{ x \in X_t \mid S_{\mathrm{ent}}(x)=1\}$.
در ادامه برای هر نمونه‌ی $x$ در $X'_t$ ، امتیاز جرم تقریبی به صورت زیر طبق رابطه‌ی \ref{eq:apms-in-batch} محاسبه می‌شود:
\begin{equation}
	\label{eq:apms-in-batch}
	S_{\mathrm{apms}}(x) = \left\| \nabla_x H(f_{\bm\theta}(x)) \right\|^2
\end{equation}
سپس وزن متناظر با هر نمونه به شکل نرمال‌سازی‌شده به صورت زیر طبق رابطه‌ی \ref{eq:apms-weights} به دست می‌آید:
\begin{equation}
	\label{eq:apms-weights}
	W(x) =
	\dfrac{
		\exp\!\left(-S_{\mathrm{apms}}(x)/T\right)
	}{
		\sum\limits_{x_i \in X'_t}
		\exp\!\left(-S_{\mathrm{apms}}(x_i)/T\right)
	}
\end{equation}

در نهایت، تابع هزینه‌ی نهایی با لحاظ کردن فیلتر $S_{\mathrm{ent}}$ و وزن‌های به‌دست‌آمده  به صورت زیر طبق رابطه‌ی \ref{eq:ocetta-objective} تعریف می‌شود:
\begin{equation}
	\label{eq:ocetta-objective}
	\mathcal{L}_{\mathrm{OCETTA}}(X'_t;\bm\theta) = \sum_{x \in X'_t} W(x)\cdot H(f_{\bm\theta}(x)),
\end{equation}

که در آن $H(\cdot)$ آنتروپی خروجی مدل است. فرآیند کلی روش پیشنهادی در قالب الگوریتم \ref{alg:ocetta} نمایش داده شده است.

\begin{algorithm}[t]
	\caption{‌انطباق زمان آزمون مستمر مبتنی بر آنتروپی با آگاهی از داده‌های خارج از توزیع}
	\label{alg:ocetta}
	\begin{latin}
		\begin{algorithmic}[1]
			\Require Stream of test samples $\mathcal{D}^{\mathrm{test}}=\{\rvx_i\}_{i=1}^{N^{\mathrm{test}}}$, 
			model $f_{\bm\theta}(\cdot)$ with trainable parameters $\tilde{\bm\theta}\subset\bm\theta$, 
			step size $\eta>0$, entropy threshold $E_0>0$.
			\Ensure Predictions $\{\hat{y}_i\}_{i=1}^{N^{\mathrm{test}}}$.
			
			\For{each online mini-batch $X_t=\{x_1,\dots,x_m\}$ from $\mathcal{D}^{\mathrm{test}}$}
			\State Compute $H(f_{\bm\theta}(x))$ and predictions 
			$\hat{y}(x) = \arg\max_c f_{\bm\theta}(x)_c$ for all $x \in X_t$
			
			\State Compute $S_{\mathrm{ent}}(x) \coloneqq \mathbb{I}\{H(f_{\bm\theta}(x)) < E_0\}$ for all $x \in X_t$
			\Comment{(Eq.~\ref{eq:hentfiltering})}
			
			\State Obtain $X'_t = \{ x \in X_t \mid S_{\mathrm{ent}}(x)=1\}$ 
			
			\State Compute $S_{\mathrm{apms}}(x) = \|\nabla_x H(f_{\bm\theta}(x))\|_2$ for all $x \in X'_t$ 
			\Comment{(Eq.~\ref{eq:apms-in-batch})}
			
			\State Compute $W(x)=\frac{\exp \!\left(-S_{\mathrm{apms}}(x)/T\right)}{\sum_{x_i \in X'_t}\exp \!\left(-S_{\mathrm{apms}}(x_i)/T\right)}$ for all $x \in X'_t$
			\Comment{(Eq.~\ref{eq:apms-weights})}
			
			\State Compute $\mathcal{L}_{\mathrm{OCETTA}}(X'_t;\bm\theta) = \sum_{x \in X'_t} W(x)\cdot H(f_{\bm\theta}(x))$ 
			\Comment{(Eq.~\ref{eq:ocetta-objective})}
			
			\State Compute gradient: $\nabla_{\tilde{\bm\theta}} \mathcal{L}_{\mathrm{OCETTA}}(X'_t;\bm\theta)$
			
			\State Update: $\tilde{\bm\theta} \leftarrow \tilde{\bm\theta} - \eta \,\nabla_{\tilde{\bm\theta}} \mathcal{L}_{\mathrm{OCETTA}}(X'_t;\bm\theta)$
			\EndFor
			
		\end{algorithmic}
	\end{latin}
\end{algorithm}



این چارچوب ترکیبی چند مزیت کلیدی دارد:
\begin{itemize}
	\item \textbf{پایداری:} با حذف یا تعدیل نمونه‌های مستعد خطا ناشی از آنتروپی زیاد یا چگالی کم ، از فروپاشی مدل در فرآیند انطباق جلوگیری می‌شود و مدل قادر است به‌صورت پایدار در جریان داده‌های زمان واقعی عمل کند.  
	\item \textbf{انعطاف‌پذیری:} برخلاف روش‌های  مبتنی بر منظم سازی مدل و تعدیل به‌روز‌رسانی وزن‌های شبکه، این روش به مدل اجازه می‌دهد دانش قبلی را حفظ کرده و در عین حال با داده‌های جدید به‌صورت پویا سازگار شود. این ویژگی، انعطاف‌پذیری الگوریتم را در شرایط تغییر توزیع تضمین می‌کند.  
	\item \textbf{کارآمدی محاسباتی:} استفاده از امتیاز جرم تقریبی و وزن‌دهی مبتنی بر اندازه‌ی گرادیان، روشی سبک و سریع ارائه می‌کند که نیازی به تنظیم پیچیده‌ی فراپارامتر‌ها ندارد و در محیط‌های عملیاتی زمان واقعی قابل استفاده است.  
	\item \textbf{تمرکز بر داده‌های با اطمینان بالا:} این ویژگی باعث می‌شود فرآیند انطباق هدفمند و بهینه باشد و دقت مدل در مواجهه با داده‌های جدید و ناشناخته افزایش یابد.  
	\item \textbf{مقاومت در برابر خطای تجمعی:} با اولویت‌دهی به نمونه‌های مطمئن و مدیریت کنترل‌شده‌ی نمونه‌های مستعد خطا، انباشت خطا و اثرات منفی داده‌های نامطمئن به حداقل می‌رسد و مدل در طول زمان عملکرد پایداری خواهد داشت.
\end{itemize}

معماری کلی روش پیشنهادی در شکل \ref{fig:ocetta} آورده شده‌است.

\begin{figure}[t]
	\centering
	\includegraphics[width=1\linewidth]{images/ocetta.pdf}
	\caption[معماری کلی روش پیشنهادی]{معماری کلی روش پیشنهادی}
	\label{fig:ocetta}
\end{figure}

به طور خلاصه، روش پیشنهادی با ترکیب مرحله‌ی پالایش نمونه‌ها و بازوزن‌دهی هوشمند مبتنی بر امتیاز جرم تقریبی، مدلی پایدار، منعطف و کارآمد را در محیط‌های پویا و جریان داده‌های زمان آزمون فراهم می‌آورد. این چارچوب امکان انطباق هوشمندانه بدون نیاز به داده‌های اولیه یا بازآموزی سنگین را فراهم می‌کند و با تمرکز بر نمونه‌های با اطمینان بالا، فرآیند یادگیری برخط را بهینه می‌سازد. 

\قسمت{شناسایی لایه‌ی بهینه در تفکیک داده‌های درون توزیع و خارج از توزیع}

در بخش قبل، روش پیشنهادی مبتنی بر امتیاز جرم تقریبی معرفی شد که بر اساس گرادیان تابع آنتروپی نسبت به ورودی $x$ عمل می‌کند و توانایی تمایز نمونه‌های درون‌توزیع و خارج از توزیع را فراهم می‌سازد. با این حال، تحقیقات اخیر نشان داده‌اند که لایه‌های مختلف شبکه‌های عصبی در شناسایی داده‌های خارج از توزیع توانایی‌های متفاوتی دارند. به طور خاص \cite{eood} نشان می‌دهد خروجی برخی لایه‌های میانی شبکه قدرت بیشتری در تفکیک داده‌های درون‌توزیع و خارج از توزیع نسبت به لایه‌های انتهایی دارند.

این مشاهده انگیزه‌ای ایجاد می‌کند تا ابتدا لایه‌ای که بیشترین تمایز را بین داده‌های درون‌توزیع و خارج از توزیع ارائه می‌دهد، شناسایی کنیم و سپس امتیاز جرم تقریبی را روی خروجی آن لایه اعمال نماییم. برای این منظور در هر لایه معیار امتیاز جرم تقریبی را برای بخشی از داده‌های درون توزیع که مدل روی آن‌ها آموزش دیده‌است و همین‌طور داده‌های خارج از توزیع اندازه‌گیری می‌کنیم و در انتهای این فرایند لایه‌ای که بیشترین تفکیک را ارائه می‌دهد انتخاب کرده و در مرحله‌ی انطباق زمان آزمون معیار امتیاز جرم تقریبی را نسبت به خروجی لایه‌ی انتخاب شده محاسبه می‌کنیم. 

در عمل، دسترسی به نمونه‌های خارج از توزیع پیش از استقرار مدل در دنیای واقعی معمولاً امکان‌پذیر نیست. بنابراین، از داده‌های درون‌توزیع موجود برای تولید نمونه‌های خارج از توزیع استفاده می‌کنیم. در این روش، از یک رویکرد پازل \lr{Jigsaw}\cite{jigsaw} استفاده می‌کنیم که تصاویر درون توزیع را به قطعه‌های $4 \times 4$ تقسیم کرده و سپس آن‌ها را به‌صورت درهم‌ریخته بازآرایی می‌کند تا تصویر خارج از توزیع نهایی ساخته شود. این فرآیند برخی ویژگی‌های نمونه‌های درون توزیع را حفظ می‌کند، اما تغییرات معنایی ایجاد می‌کند که شناسایی آن‌ها را به چالش می‌کشد. اگر مجموعه داده‌ی آموزشی درون توزیع را به‌صورت $X = \{x_i\}_{i=1}^I$ داشته باشیم، مجموعه‌ی خارج از توزیع تولید شده با $\hat{X} = \{\hat{x}_i\}_{i=1}^I$ نشان داده می‌شود. نمونه‌ای از اعمال این تبدیل در شکل \ref{fig:jigsaw} آورده شده‌است.

\begin{figure}[t]
	\centering
	\subfigure{	
		\includegraphics[width=0.23\linewidth]{images/jigsaw1.png}
	}
	\subfigure{
		\includegraphics[width=0.23\linewidth]{images/jigsaw2.png}
	}
	\subfigure{
		\includegraphics[width=0.23\linewidth]{images/jigsaw3.png}
	}
	\subfigure{
		\includegraphics[width=0.23\linewidth]{images/jigsaw4.png}
	}
	\caption[تبدیل \lr{Jigsaw}]{ نمونه‌ای از اعمال تبدیل \lr{Jigsaw} که تصاویر درون‌توزیع را برای تولید نمونه‌های خارج از توزیع مصنوعی به قطعات کوچکتر تقسیم کرده و بازآرایی می‌کند.}
	\label{fig:jigsaw}
\end{figure}

برای هر لایه‌ی کاندید $l \in L$، امتیاز جرم تقریبی نسبت به خروجی لایه طبق رابطه‌ی \ref{eq:apms_layer} تعریف می‌شود. اگر $h_l(x)$ خروجی لایه $l$ برای نمونه $x$ باشد، داریم:  
\begin{equation}
	\label{eq:apms_layer}
	S_{\mathrm{apms}}(x,l) = \left\| \nabla_{h_l} H(f_{\bm\theta}(x)) \right\|_2,
\end{equation}  
که در آن $p_\theta(x)$ چگالی مدل و $\nabla_{h_l}$ گرادیان نسبت به خروجی لایه $l$ را نشان می‌دهد. این تعریف امکان ارزیابی حساسیت هر لایه نسبت به تغییرات اطلاعات توزیع را فراهم می‌آورد.  

برای هر نمونه $x \in X$ و نمونه‌ی خارج از توزیع تولید شده‌ی $x' \in X'$ امتیاز $S_{\mathrm{apms}}(x,l)$ و $S_{\mathrm{apms}}(x',l)$ محاسبه شده و در مجموعه‌های $S_{\mathrm{id}}(l)$ و $S_{\mathrm{ood}}(l)$ ذخیره می‌شوند.  

پس از محاسبه‌ی امتیازهای لایه‌ای برای تمام نمونه‌ها، برای هر لایه یک معیار ترکیبی $S(l)$ به صورت نرمال‌سازی اختلاف بین نمونه‌های درون توزیع و نمونه‌های خارج از توزیع طبق رابطه‌ی \ref{eq:layer_score} تعریف می‌شود:  
\begin{equation}
	\label{eq:layer_score}
	S(l) = -\frac{1}{I} \sum_{i=1}^{I} \frac{S_{\mathrm{id}}^{(i)}(l) - S_{\mathrm{ood}}^{(i)}(l)}{\max_{1\le j \le I} \left| S_{\mathrm{id}}^{(j)}(l) - S_{\mathrm{ood}}^{(j)}(l) \right|}.
\end{equation}  
لایه‌ی با بیشترین مقدار $S(l)$ طبق رابطه‌ی \ref{eq:layer_argmax} به عنوان لایه‌ی حساس $\ell^*$ انتخاب می‌شود:  
\begin{equation}
	\label{eq:layer_argmax}
	\ell^* = \arg\max_{l \in L} S(l).
\end{equation}  

انتخاب لایه‌ی مناسب امکان اعمال مستقیم امتیاز جرم تقریبی روی فضای ویژگی‌های میانی شبکه را فراهم می‌کند و بدین ترتیب فرآیند شناسایی داده‌های خارج از توزیع با دقت و پایداری بیشتری انجام می‌شود. روند کامل این فرآیند در قالب الگوریتم \ref{alg:layer_selection} ارائه شده است.  

\begin{algorithm}[t]
	\caption{شناسایی لایه‌ی بهینه در تفکیک داده‌های درون توزیع و خارج از توزیع}
	\label{alg:layer_selection}
	\begin{latin}
		\begin{algorithmic}[1]
			\Require Training data $\mathcal{D}^{\mathrm{train}} = \{x_i\}_{i=1}^I$, 
			candidate layers $L = \{l_1,\dots,l_K\}$
			\Ensure Selected layer index $\ell^*$ (best separator between ID and OOD)
			
			\State Initialize $S_{\mathrm{id}}(l) \gets \varnothing,\; S_{\mathrm{ood}}(l) \gets \varnothing \quad \forall l \in L$
			\State Initialize $S(l) \gets 0 \quad \forall l \in L$
			
			\For{each mini-batch $X_t = \{x_1,\dots,x_m\}$ from $\mathcal{D}^{\mathrm{train}}$}
			\State Obtain $X'_t = \{\mathrm{Jigsaw}(x) \mid x \in X_t\}$
			\For{each $l \in L$}
			\State Compute $S_{\mathrm{apms}}(x,l)$ for all $x \in X_t$, append them into $S_{\mathrm{id}}(l)$
			\Comment{(Eq.~\ref{eq:apms_layer})}
			\State Compute $S_{\mathrm{apms}}(x',l)$ for all $x' \in X'_t$, append them into $S_{\mathrm{ood}}(l)$
			\Comment{(Eq.~\ref{eq:apms_layer})}
			\EndFor
			\EndFor
			
			\For{each layer $l \in L$}
			\State Compute $S(l) = -\frac{1}{I} \sum_{i=1}^{I} 
			\frac{S_{\mathrm{id}}^{(i)}(l) - S_{\mathrm{ood}}^{(i)}(l)}
			{\max_{1 \le j \le I} \left| S_{\mathrm{id}}^{(j)}(l) - S_{\mathrm{ood}}^{(j)}(l) \right|}$
			\Comment{(Eq.~\ref{eq:layer_score})}
			\EndFor
			
			
			\State $\ell^* \gets \arg\max_{l \in L} S(l)$
			\Comment{(Eq.~\ref{eq:layer_argmax})}
			\State \textbf{return} $\ell^*$
		\end{algorithmic}
	\end{latin}
\end{algorithm}

\قسمت{جمع‌بندی}

در این فصل، رویکرد پیشنهادی برای انطباق در زمان آزمون برخط ارائه گردید. نقطه‌ی آغاز بحث، شناسایی یکی از مهم‌ترین ضعف‌های شبکه‌های عصبی عمیق، یعنی پدیده‌ی اطمینان بیش‌ازحد بود. همان‌طور که در بخش‌های ابتدایی فصل توضیح داده شد، این پدیده باعث می‌شود که شبکه‌های عصبی در مواجهه با داده‌های خارج از توزیع نیز پیش‌بینی‌هایی با سطح اطمینان بالا ارائه دهند، حتی زمانی که پاسخ به‌طور کامل نادرست باشد. این رفتار نه‌تنها قابلیت اعتماد مدل را کاهش می‌دهد، بلکه در حالت انطباق برخط در زمان آزمون که هیچ برچسبی در دسترس نیست، می‌تواند مدل را به سمت انطباق نادرست و انباشت خطا سوق دهد. از این‌رو، طراحی روشی که در برابر این پدیده مقاوم باشد و بتواند سطح اطمینان مدل را به شکل هوشمندانه‌ای کنترل کند، یک ضرورت محسوب می‌شود.  

با در نظر گرفتن این چالش، ابتدا یک مرور تحلیلی بر رویکردهای پیشین در حوزه‌ی شناسایی داده‌های خارج از توزیع انجام شد. نشان داده شد که بسیاری از روش‌های آموزشی نیازمند داده‌های کمکی یا بازآموزی کامل شبکه هستند و بنابراین در حالتهای عملی، خصوصاً در محیط‌های برخط و پویا، قابلیت استفاده‌ی محدودی دارند. همچنین روش‌های پسینی که بدون نیاز به بازآموزی ارائه شده‌اند، عمدتاً بر امتیازهایی مانند آنتروپی یا بیشینه‌ی خروجی شبکه تکیه دارند؛ در حالی‌که این امتیازها در حضور داده‌های مخدوش یا خارج از توزیع می‌توانند گمراه‌کننده باشند.  

برای رفع این محدودیت‌ها، معیار جدیدی با عنوان امتیاز جرم تقریبی معرفی شد. این معیار بر مبنای اندازه‌ی گرادیان لگاریتم درست‌نمایی نسبت به ورودی تعریف گردید و مزیت اصلی آن در نظر گرفتن تغییرات محلی چگالی پیرامون هر نمونه است. بدین ترتیب، برخلاف روش‌های مبتنی بر چگالی نقطه‌ای، این معیار قادر است داده‌هایی را که به‌رغم داشتن مقدار درست‌نمایی بالا، در نواحی کم‌جرم و غیرنمونه‌وار قرار دارند، شناسایی کند. این ویژگی کلیدی سبب می‌شود که امتیاز جرم تقریبی ابزاری کارآمد برای تمایز میان نمونه‌های درون‌توزیع و خارج از توزیع در جریان داده‌های برخط باشد.  

در ادامه، سازوکاری دو مرحله‌ای برای بهره‌گیری از این معیار در فرایند انطباق ارائه گردید. در گام نخست، داده‌های ورودی بر اساس آنتروپی پالایش شدند تا نمونه‌های آشکارا نامطمئن حذف گردند. سپس، در گام دوم، داده‌های باقی‌مانده با توجه به امتیاز جرم تقریبی وزن‌دهی شدند. بدین‌ترتیب، نمونه‌های قابل اعتماد سهم بیشتری در به‌روزرسانی مدل دریافت کرده و داده‌های مشکوک یا خارج از توزیع اثر کمتری بر انطباق داشتند. این طراحی می‌تواند به شکل مستقیم به کاهش خطر انباشت خطا و افزایش پایداری مدل در طول زمان منجر شود.  

چارچوب پیشنهادی به‌طور ویژه متناسب با شرایطی در نظر گرفته شده است که در آن برچسب داده‌های ورودی در زمان آزمون در دسترس نیست، امکان تنظیم دقیق فراپارامترها وجود ندارد، منابع محاسباتی محدود هستند و پایداری انطباق در حضور داده‌های خارج از توزیع اهمیت حیاتی دارد. رویکرد پیشنهادی با ترکیب دو سازوکار اصلی شامل پالایش اولیه‌ی داده‌ها بر اساس آستانه‌ی آنتروپی و وزن‌دهی پویا به نمونه‌ها با بهره‌گیری از امتیاز جرم تقریبی، زمینه‌ی انطباقی پایدار و مقاوم را فراهم می‌سازد.  

ویژگی بارز این روش در مقایسه با بسیاری از رویکردهای موجود، سادگی ساختار، نیازمند نبودن به بازآموزی کامل شبکه و قابلیت به‌کارگیری مستقیم بر روی مدل‌های از پیش آموزش‌دیده است. این خصیصه‌ها آن را برای استفاده در حالتهای واقعی که محدودیت‌های زمانی و محاسباتی جدی وجود دارد، مناسب می‌سازد.  

در فصل بعد، با تکیه بر نتایج تجربی نشان داده خواهد شد که این چارچوب پیشنهادی علاوه بر ویژگی‌های مفهومی یادشده، از نظر عملکرد نیز توانایی مقابله با چالش‌های مطرح‌شده را دارد و در عمل می‌تواند دقت و پایداری مدل را در شرایط متغیر و پویا حفظ کند.